\chapter{Findings and Discussion}

\section{Introduction}
This chapter presents the findings from the thematic analysis of the selected literature and discusses their implications for understanding consumer perception in AI-driven marketing. Rather than treating positive and negative outcomes as contradictory findings, the discussion interprets them as recurring trade-offs and conditional effects. Following Puntoni et al.'s experiential framework \cite{puntoni2021experiential}, the chapter is organised around four recurring tensions: data capture, classification, delegation, and social interaction. This structure makes it possible to show not only what patterns appear in the literature, but also why those patterns arise and under what conditions consumer responses become more favourable or more resistant.

For clarity, the analysis distinguishes between three levels that are closely related but not identical: interpretive appraisal, affective response, and behavioural outcome. In the chapter, consumer perception refers most directly to the first two levels, that is, how consumers make sense of AI and how they feel about it. Behavioural outcomes such as engagement, purchase intention, avoidance, or resistance are treated as consequences shaped by those perceptions rather than as perception itself.

\section{Interpreting the Four Tensions}
The four thematic groups identified in this dissertation should not be understood as four separate lists of positive and negative findings. Rather, each of them captures a recurring tension in how consumers interpret AI in marketing. These tensions help explain why consumer perception of AI is rarely stable or one-dimensional. Across the literature, consumers often value what AI provides while simultaneously resisting the conditions through which those benefits are delivered.

The first tension is between being served and being exploited. AI-driven marketing can make the consumer experience more relevant, efficient, and convenient, but these benefits depend on extensive data capture. As a result, consumers are often drawn into a privacy calculus in which functional value is weighed against surveillance, loss of control, and mistrust. The second tension is between being understood and being misunderstood. AI classification and personalisation can make consumers feel recognised, yet the same systems may also make them feel stereotyped, reduced to a category, or stripped of individuality.

The third tension is between being empowered and being replaced. Delegating tasks to AI can reduce effort and improve efficiency, but it can also undermine autonomy, self-efficacy, and the sense of authorship over one's own choices. The fourth tension is between being connected and being alienated. AI-mediated communication can create speed, responsiveness, and even novelty, but it may also feel emotionally empty or morally inappropriate when consumers expect genuine empathy, warmth, or human sincerity.

Taken together, these tensions suggest that the literature does not simply contain contradictory findings. Instead, it shows that consumer responses to AI are shaped by recurring trade-offs. AI is more likely to be perceived positively when it delivers practical value without violating expectations of privacy, individuality, autonomy, or authenticity. It is more likely to be perceived negatively when those benefits come at too high a psychological, emotional, or ethical cost. The four sections that follow develop this argument in greater detail.

\section{Overview of the Coding Framework}
The review suggests that consumer perception of AI is shaped by four recurring tensions:

\begin{itemize}
    \item Data capture: served versus exploited
    \item Classification: understood versus misunderstood
    \item Delegation: empowered versus replaced
    \item Social interaction: connected versus alienated
\end{itemize}

These tensions do not operate independently. Across the literature, consumer responses to AI are shaped by the interaction between perceived benefits such as convenience, relevance, efficiency, and responsiveness, and perceived costs such as intrusiveness, opacity, stereotyping, loss of autonomy, and perceived inauthenticity. The findings therefore suggest that AI in marketing should be understood through a relational and interpretive lens, not simply through a performance lens. Trust, authenticity, privacy concern, and behavioural intention are best understood as outcomes or mechanisms distributed across these four broader experience domains rather than as fully separate themes.

\section{Data Capture and Privacy: Served versus Exploited}
The first thematic tension concerns data capture. AI-driven marketing depends heavily on consumer data in order to personalise recommendations, optimise advertising, and support real-time interaction. From the consumer perspective, this can create a positive experience when data use produces relevance, convenience, and reduced search effort. In such cases, consumers may feel served by AI because the system appears to anticipate their needs and improve the efficiency of the market experience \cite{an2025personalized,ameen2021customer}.

However, the same process can produce a very different perception when data collection feels opaque, excessive, or difficult to control. Aguirre et al. show that personalised advertising performs better when information collection is overt, but that covert collection heightens perceived vulnerability and reduces effectiveness \cite{aguirre2015paradox}. Hardcastle et al. similarly show that AI-driven personalised journeys can create frustration and resentment when consumers feel persistently monitored or when stored data seems overly extensive \cite{hardcastle2025journeys}. Boerman et al. likewise show that targeted digital advertising can quickly shift from relevance to perceived intrusiveness when consumers become more aware of how data is being collected and used \cite{boerman2017oba}. Martin and Murphy provide broader conceptual support by showing that privacy in marketing is not a peripheral issue but a structural concern tied to trust, legitimacy, and power in data use \cite{martin2017privacy}. More broadly, Davenport et al. and Kumar et al. identify privacy and data governance as structural tensions in AI-enabled marketing \cite{davenport2020future,kumar2024aipowered}. Ozturkcan and Bozda{\u{g}} extend this concern by showing how opaque practices and exaggerated claims about AI can generate a broader cycle of mistrust, in which AI washing is followed by consumer backlash or ``AI booing'' \cite{ozturkcan2025mistrust}. Srivastava and Gurme further show that privacy concern and perceived surveillance can directly weaken trust even when personalisation remains technically effective \cite{srivastava2026social}. These studies suggest that the issue is not simply whether firms use data, but whether that use is interpreted as proportionate, transparent, and legitimate.

What consumers perceive in this theme is a tension between being efficiently served and being quietly exploited. Consumers value tailored relevance, convenience, and reduced search effort, but they also become highly sensitive to signs that their personal data is being continuously monitored, stored, or abstracted into computational profiles without meaningful consent. This is why the same AI system can be interpreted either as useful service or as invasive surveillance.

Why this perception occurs can be explained through privacy calculus and perceived vulnerability. Consumers weigh the practical benefits of personalisation against the costs of dataveillance, loss of privacy, and reduced control. This tension becomes especially intense when data collection is covert or when the logic of the system remains opaque. It is further exacerbated when firms overstate their AI sophistication or ethical responsibility, because AI washing inflates expectations that are then violated when the system appears exploitative rather than trustworthy. The underlying mechanism is therefore not merely data use itself, but the interpretation of that data use as either legitimate service or covert extraction.

The main behavioural outcomes affected by this tension are trust, click-through intention, ad avoidance, purchase intention, and longer-term willingness to remain engaged with the customer journey. Transparent and proportionate data practices are more likely to sustain trust and purchase intention, whereas opaque or excessive data capture is more likely to produce avoidance, reactance, and, in more severe cases, public backlash or disengagement \cite{an2025personalized,srivastava2026social}. The evidence for this theme is especially strong because similar patterns appear across experiments, surveys, and review-based studies. The tension is therefore best understood not as a contradiction in the literature, but as a privacy-personalisation trade-off built into the data capture experience itself.

\section{Classification and Personalisation: Understood versus Misunderstood}
The second thematic tension concerns classification. AI systems sort consumers into categories, infer their likely preferences, and use those classifications to generate personalisation. When this process works well, consumers may feel understood because the message, recommendation, or service appears relevant to their actual needs. This is why personalisation is often associated with favourable consumer response in the literature \cite{an2025personalized,haleem2022ai,kumar2024aipowered}.

Yet the classification process can also make consumers feel misunderstood. Puntoni et al. suggest that AI can produce experiences of misunderstanding when it reduces people to narrow inferred categories or makes judgements that fail to reflect their full identity and intention \cite{puntoni2021experiential}. Hou et al. provide particularly direct evidence for this problem by showing that consumers often infer that AI-generated personalised solutions are shared with a larger group of people and are therefore less unique \cite{hou2026uniqueness}. Starke et al. further show that perceptions of algorithmic fairness are highly sensitive to context, procedure, and outcome, which helps explain why classification can also trigger concerns about stereotyping, bias, and unequal treatment \cite{starke2022fairness}. Longoni et al. add an important boundary condition by showing that consumers become especially resistant when AI is perceived to lack the judgement needed for high-stakes or personally nuanced decisions \cite{longoni2019medical}. Even where personalisation is technically accurate, it may still feel mechanical or impersonal if the consumer believes that many others are being treated in essentially the same way. The issue is therefore not only predictive success, but whether the classification feels sufficiently individual and context-sensitive.

What consumers perceive in this theme is a tension between being recognised and being reduced. On the positive side, AI classification can make consumers feel understood when the system appears to anticipate needs and provide relevant solutions. On the negative side, consumers may feel misunderstood when they suspect that the system relies on stereotypes, rigid templates, or overgeneralised profiles rather than on an appreciation of their individual complexity.

Why this perception occurs can be explained through uniqueness neglect and fairness concerns. Because consumers often attribute low experiential sensitivity to AI, they assume that algorithmic recommendations are generated through broad rules rather than genuine insight. This leads them to infer a larger expected group size, meaning that the personalised solution does not feel rare, exclusive, or truly individual. Where biased or historically skewed data are involved, this perception becomes more severe because classification can also be interpreted as unfair or discriminatory. The central mechanism is therefore not simply whether the recommendation is personalised, but whether the underlying classification appears psychologically sensitive and procedurally fair.

The behavioural outcomes affected by this theme include adoption intention, willingness to rely on AI advice, willingness to pay for personalised solutions, and acceptance of AI providers in higher-stakes contexts. When classification produces a convincing sense of understanding, AI is more likely to be accepted as useful and competent. When it produces uniqueness neglect or fairness concerns, consumers become more reluctant to trust the recommendation, less willing to adopt it, and more likely to reject AI providers in contexts where individual nuance matters. The evidence base here is moderate to strong: the overall pattern is clear, even though fairness perceptions remain highly context-sensitive. The classification experience therefore introduces a subtle but important trade-off between precision and reductionism.

\section{Task Delegation and Autonomy: Empowered versus Replaced}
The third thematic tension concerns delegation. Consumers often rely on AI to assist with search, recommendation, customer service, and decision support. In many contexts this delegation is experienced positively because it reduces effort, saves time, and makes the market environment easier to navigate. Ameen et al. show that AI-enabled services can improve customer experience when they are perceived as competent, convenient, and supportive of the consumer's goals \cite{ameen2021customer}. In these situations, delegation is experienced as empowerment rather than dependence.

At the same time, delegation can generate a loss of autonomy. Puntoni et al. argue that AI may create tension when consumers feel that judgement, choice, or effort has been transferred too far away from the self \cite{puntoni2021experiential}. Hardcastle et al. suggest that this tension intensifies when predictive journeys repeatedly direct consumers toward similar content or options, effectively narrowing what they see and reinforcing a sense of constrained choice \cite{hardcastle2025journeys}. Davenport et al. similarly note that AI-driven micro-targeting can reduce search costs while also undermining the consumer's perceived freedom of choice \cite{davenport2020future}. This becomes especially relevant when AI does not merely support a decision but appears to narrow the range of options, shape preferences too strongly, or replace forms of action that consumers associate with agency, creativity, or competence. In such cases, the consumer no longer experiences AI as assistance but as partial replacement.

What consumers perceive in this theme is a tension between empowerment and displacement. AI can feel empowering when it removes mundane effort and helps consumers navigate large amounts of information more efficiently. Yet the same delegation can feel threatening when consumers sense that control is shifting away from them, especially in tasks that involve subjectivity, judgement, or identity.

Why this perception occurs can be explained through perceived sacrifice and loss of autonomy. AI systems can guide choices so heavily that consumers no longer experience themselves as the true authors of their decisions. This becomes more serious when algorithmic curation narrows visible options or creates filter bubbles, because the convenience of delegation starts to look like a loss of freedom. In identity-relevant activities, delegation may also feel like the loss of an accomplishment that consumers believe should belong to them rather than to a machine. The core mechanism here is the conversion of convenience into psychological threat when support turns into substitution.

The behavioural outcomes associated with this theme include resistance to recommendations, frustration, reduced reliance on AI support, and psychological reactance. In some cases, consumers may deliberately ignore or reject algorithmic guidance simply to reassert their independence. Although the empirical base on long-term autonomy loss is still developing, the existing evidence consistently suggests that delegation is most positively perceived when it remains supportive and reversible, and most negatively perceived when it appears to replace rather than augment the consumer's own agency.

\section{Social Interaction and Authenticity: Connected versus Alienated}
The fourth thematic tension concerns social interaction. Consumers increasingly encounter AI in socially framed contexts, such as chatbots, AI-generated advertising, virtual assistants, and other forms of automated brand communication. These interactions can create a sense of connection when they are responsive, clear, and useful. Consumers may appreciate instant assistance, reduced waiting time, and the novelty of AI-mediated brand interaction, especially in factual or utilitarian settings \cite{ameen2021customer,chan2025genai}.

However, this same social experience can become alienating when AI communication feels emotionally empty, morally inappropriate, or insufficiently human. This is where authenticity becomes especially important. As discussed in the literature review, authenticity in this dissertation has two dimensions: source authenticity and emotional authenticity. Source authenticity concerns whether the message is perceived as genuinely human-authored or AI-authored \cite{kirk2025aiauthorship}. Emotional authenticity concerns whether the communication feels sincere, empathetic, and genuinely intended \cite{arango2023charitable}. These two dimensions help explain why consumers may resist AI not simply because it is artificial, but because it disrupts expected forms of warmth, empathy, or symbolic meaning.

What consumers perceive in this theme is a dual response: AI can feel novel, responsive, and efficient, yet also inauthentic, emotionally empty, or even deceptive. Consumers may enjoy the speed and availability of AI-mediated communication, but they remain aware that machines do not possess genuine internal emotional states. Kirk and Givi show that this can produce an AI-authorship effect in which the perceived lack of real feeling weakens response to emotionally expressive communication \cite{kirk2025aiauthorship}. Shi and Jiang similarly show that AI disclosure can increase novelty while simultaneously reducing authenticity \cite{shi2026disclosure}. This creates a tension between the convenience of connection and the discomfort of artificiality.

Why this perception occurs can be explained through two competing mechanisms. The first is perceived novelty, which can stimulate curiosity and engagement because AI signals innovation. The second, often more powerful, is perceived inauthenticity. When AI attempts to simulate empathy, warmth, or emotional sincerity, consumers may judge the interaction as fake because the expression is not backed by genuine feeling. Arango et al. show how this problem can reduce empathy and donation intention in charitable advertising \cite{arango2023charitable}, while Belanche et al. show that resistance becomes stronger in hedonic or high-involvement settings where realism and genuine human presence matter more \cite{belanche2025images}. In some contexts this mismatch generates moral disgust, especially when AI is used in emotionally charged or morally sensitive communication. The effect is therefore strongly context-dependent: utilitarian tasks invite more acceptance, whereas hedonic or high-involvement experiences intensify the need for human touch and authentic presence.

The behavioural outcomes associated with this theme include short-term engagement through novelty, but also reduced word-of-mouth, lower loyalty, weaker donation intention, and stronger resistance when authenticity concerns dominate. In medical, charitable, and emotionally expressive settings, alienation can become especially strong because consumers fear that their individuality, seriousness, or emotional needs are not being genuinely recognised \cite{longoni2019medical}. The evidence for this theme is strong across multiple experiments and applied contexts. Social interaction with AI is therefore best understood as a double-edged mechanism that can create both connection and alienation depending on how the communication is framed and where it occurs.

\section{How Trust, Authenticity, and Behavioural Intention Fit Across the Four Themes}
The analysis suggests that trust, authenticity, and behavioural intention should not be treated as isolated findings detached from the broader experience structure. In the literature review, trust was discussed as a distinct theme because it is one of the most visible concepts in prior studies and often appears as an explicit dependent or mediating variable. In the findings, however, trust is repositioned as a cross-cutting mechanism because the thematic synthesis shows that it is produced through several different experience domains rather than operating as one separate empirical topic. Trust is shaped across data capture, classification, delegation, and social interaction, which is why it is analytically more useful here as an evaluative mechanism than as a fifth standalone theme.

Authenticity is primarily concentrated in the social interaction theme, but it also interacts with trust because doubts about authorship and emotional sincerity can weaken the credibility of AI-mediated communication. Behavioural intention, by contrast, is best understood as an outcome of how the four tensions are resolved. Consumers are more likely to engage, purchase, or respond favourably when they feel served, understood, empowered, and connected. They are less likely to do so when they feel exploited, misunderstood, replaced, or alienated.

\section{How the Four Tensions Interact}
The four tensions identified in this dissertation should not be treated as independent blocks. Although they are presented separately for analytical clarity, the literature suggests that they often reinforce or condition one another. In practice, consumers do not encounter data capture, classification, delegation, and social interaction as isolated dimensions. They encounter them together within the same digital experience, which means that one tension can intensify or moderate another.

The clearest interaction appears between data capture and classification. Consumers are less likely to feel understood by personalisation when they are already uneasy about how their data was collected. In this sense, privacy concern can weaken the positive effects of relevance, because the same recommendation may be interpreted less as helpful recognition and more as intrusive profiling. A similar interaction exists between classification and delegation. When AI appears to rely on narrow categories or generic rules, consumers are less willing to delegate judgement to it, especially in contexts that require nuance, taste, or personal identity. Misunderstanding therefore makes replacement feel more threatening.

Data capture also interacts with social interaction. AI-mediated communication may already feel emotionally thin or inauthentic, but that concern can become stronger when the message is also understood as the product of hidden surveillance or excessive profiling. In other words, perceived intrusion can make perceived inauthenticity more severe. Conversely, when AI use is transparent, proportionate, and clearly bounded, consumers may be more willing to accept novelty or limited automation in communication.

Taken together, these patterns suggest that the four tensions form an interpretive system rather than a simple list. Data concerns can shape whether personalisation feels caring or exploitative. Classification can shape whether delegation feels supportive or substitutive. Social interaction can amplify how questions of trust, privacy, and authenticity are ultimately resolved. This interactional reading strengthens the framework by showing that consumer perception is shaped not only by the presence of individual tensions, but by how those tensions accumulate and condition each other within a single AI-mediated encounter.

\section{Conditions That Shape More Positive or Negative Responses}
The literature suggests that consumer responses to AI become more positive or more negative under identifiable conditions. One important condition is the visibility of AI. When AI works in the background to improve convenience or targeting, it is often evaluated more positively. When it becomes highly visible as the creator of emotionally meaningful or identity-relevant communication, consumers become more alert to questions of authenticity and trust.

Another condition is message type. AI appears to be more favourably received when the message is informational, functional, or service-oriented than when it is expected to convey emotional sincerity, artistic originality, or social meaning. Product and context also matter. High-involvement, hedonic, and symbolic categories tend to intensify concerns about authenticity and human judgement, whereas lower-stakes or convenience-driven contexts often make efficiency benefits more salient \cite{belanche2025images,arango2023charitable}.

Data practice is a further condition. Consumers are more willing to accept AI when personalisation feels proportionate and understandable. They become more resistant when the same personalisation feels opaque, excessive, or difficult to control \cite{boerman2017oba,srivastava2026social}. These patterns suggest that the consumer response to AI is not random. It is shaped by recurring contextual conditions that determine whether the benefits of AI outweigh its relational and ethical costs.

\section{Implications of the Findings}
For researchers, the findings suggest that future work should move beyond simply identifying whether AI has positive or negative effects. What matters more is specifying the mechanisms and conditions through which those effects occur. This would allow future studies to build a more cumulative understanding of consumer response to AI and GenAI across different marketing settings.

For practitioners, the main implication is that AI should not be managed only as an efficiency tool. It must also be managed as a communicative and relational presence. The success of AI in marketing depends on whether it is used in ways that preserve trust, respect privacy expectations, and match the symbolic demands of the message context. In practice, this means that high precision, fast content generation, or automated interaction are not enough on their own. They must be aligned with what consumers interpret as appropriate, credible, and worth engaging with.

\section{Summary}
This chapter has shown that consumer perception of AI in marketing is shaped by recurring trade-offs rather than by simple contradiction. Organised through the four tensions of data capture, classification, delegation, and social interaction, the analysis suggests that AI tends to generate more positive responses when consumers feel served, understood, empowered, and connected, but more negative responses when they feel exploited, misunderstood, replaced, or alienated. These outcomes are produced through identifiable mechanisms and shaped by identifiable conditions, including AI visibility, message type, emotional expectation, and data practice. The next chapter concludes the dissertation by summarising these findings and outlining their broader implications.
